\chapter{Trimmed surfaces}\label{ch:trimming}

A NURBS surface is always a deformed \emph{rectangle}: its parameter domain is
$[u_0,u_1]\times[v_0,v_1]$. Real faces are rarely rectangular -- a plate has a bore, a lens
has a circular aperture, the side of a drilled cylinder has an elliptic window. \NF{}
represents such faces as \emph{trimmed surfaces}: a NURBS surface together with closed loops
of curves in its parameter domain that cut out the part that belongs to the face.

This chapter explains the trimming concept, the loop and face types
\apiref{TrimLoop} and \apiref{TrimmedSurface}, and the two \apiref{split} functions that
cut surfaces into trimmed regions. Trimmed surfaces are the faces of shells and solids
(\cref{ch:solids}), they are produced by Booleans, and they are transferred exactly to gmsh
(\cref{ch:meshing}) and STEP (\cref{ch:step}).

\section{Trimming in parameter space}\label{sec:trim-concept}

The loops of a trimmed surface live in the $(u,v)$ parameter plane of the surface, not in 3D.
A trim curve is an ordinary \code{NurbsCurve} whose points are $(u, v, 0)$: the $x$ coordinate
is $u$, the $y$ coordinate $v$, and $z$ is zero. The face is the image of the region bounded
by the loops:
\[
F = \{\, S(u,v) \;:\; (u,v) \text{ inside the outer loop and outside all holes} \,\}.
\]

This representation has three important consequences:

\begin{itemize}
\item \textbf{It is exact.} The trimmed face is exactly the part of the exact surface;
nothing is tessellated. A circle in the parameter plane of a flat rectangle maps to an exact
circle or ellipse in 3D.
\item \textbf{Shapes are distorted by the parametrisation.} A circle of radius $0.25$ in
$(u,v)$ on a $2\times2$ square ($u = x/2$, $v=y/2$) is a circle of radius $0.5$ in 3D; on a
$4\times2$ plate it is an ellipse with semi-axes $1$ and $0.5$. On curved surfaces the image
can be any curve. Always think about the map $(u,v)\mapsto S(u,v)$ when you build loops by
hand.
\item \textbf{The 3D boundary is derived.} The 3D edges of a face are computed by composing
the trim curves with the surface (\code{boundary_curves}, \cref{sec:trim-boundary}); when faces
are sewn into shells these images are matched to find shared edges.
\end{itemize}

\subsection{Orientation rules}

Loops are oriented in the $(u,v)$ plane:

\begin{itemize}
\item the \textbf{outer loop} runs \textbf{counter-clockwise} (positive signed area);
\item \textbf{holes} run \textbf{clockwise} (negative signed area);
\item hence the face always lies \emph{to the left} of every loop, seen in the $(u,v)$ plane
from $+z$.
\end{itemize}

You do not have to get this right yourself: the \code{TrimmedSurface} constructor orients
the loops on construction (reversing those that run the wrong way). The rule matters when you
read loops back -- \code{face.outer}, \code{face.holes} -- or when you interpret section loops
(\cref{ch:slicing}), which follow the same convention.

The surface orientation, i.e. the direction of the normal $S_u\times S_v$, is a different
matter: it is a property of the surface parametrisation, not of the loops. It is changed by
\code{flipped()} (\cref{sec:trim-flip}).

\section{Trim loops}\label{sec:trimloop}

A \apiref{TrimLoop} is a closed chain of parameter-space curves: the end of each curve is
the start of the next, and the last curve ends at the start of the first.

\begin{lstlisting}
loop = nf.TrimLoop([c0, c1, c2])   # curves with z = 0; x = u, y = v
\end{lstlisting}

Three helper functions create the common loops:

\begin{center}
\begin{tabularx}{\linewidth}{@{}lX@{}}
\toprule
Function & Loop \\
\midrule
\apiref{rectangle_loop}\code{(u0, u1, v0, v1)} & rectangle $[u_0,u_1]\times[v_0,v_1]$, four
lines, counter-clockwise \\
\apiref{circle_loop}\code{(cu, cv, r)} & \emph{exact} circle (rational quadratic NURBS) of
radius $r$ about $(c_u, c_v)$, counter-clockwise \\
\apiref{polygon_loop}\code{(points)} & closed polygon through the given $(u,v)$ points
(the closing segment is added) \\
\bottomrule
\end{tabularx}
\end{center}

A \code{TrimLoop} has the following members:

\begin{center}
\begin{tabularx}{\linewidth}{@{}lX@{}}
\toprule
Member & Meaning \\
\midrule
\code{curves} & the curves of the loop, in order \\
\code{signed_area()} & signed area enclosed in the parameter plane: $>0$ for
counter-clockwise, $<0$ for clockwise loops \\
\code{reversed()} & the same loop traversed in the opposite direction \\
\code{polyline(samples_per_span=16)} & dense polyline approximation as a list of $[u, v]$
pairs (the first point is not repeated at the end) \\
\bottomrule
\end{tabularx}
\end{center}

\begin{note}
\code{signed_area()} is computed from the dense polyline, so for curved loops it is accurate
to the polyline resolution only: the circle loop of radius $0.25$ in the example below reports
$0.196032$, while the exact area $\pi/16$ is $0.196350$. It is intended for orientation tests,
not for measurements. Use \code{TrimmedSurface.area()} for exact areas.
\end{note}

\section{Trimmed surfaces}\label{sec:trimmed-surface}

\apiref{TrimmedSurface} has two constructors:

\begin{lstlisting}
nf.TrimmedSurface(surface)                                     # untrimmed: outer loop = domain
nf.TrimmedSurface(surface, outer, holes=[], tolerance=nf.Tolerance())
\end{lstlisting}

The loops must be closed (within \code{tolerance.length}, measured in parameter units) and
must lie in the parameter domain; otherwise \code{GeometryError} is raised. The holes must
lie inside the outer loop.

\begin{center}
\begin{tabularx}{\linewidth}{@{}lX@{}}
\toprule
Member & Meaning \\
\midrule
\code{surface} & the underlying \code{NurbsSurface} \\
\code{outer}, \code{holes} & the loops (oriented: outer counter-clockwise, holes clockwise) \\
\code{is_untrimmed} & \code{True} if the region is the whole parameter domain \\
\code{contains(u, v)} & \code{True} if $(u,v)$ lies inside the region (on the boundary: either
answer) \\
\code{area(rel_tol=1e-9)} & area of the face in 3D \\
\code{boundary_curves(tolerance=1e-7)} & the loops mapped to 3D, outer loop first \\
\code{with_hole(loop)} & a new face with an additional hole \\
\code{flipped()} & the same face with opposite normal \\
\code{label}, \code{with_label(label)} & free-form name of the face \\
\bottomrule
\end{tabularx}
\end{center}

Like all \NF{} geometry, trimmed surfaces are immutable: \code{with_hole},
\code{flipped} and \code{with_label} return new objects.

\subsection{A worked example}

\example{10_trimmed_faces.py}

\begin{lstlisting}[style=text]
rectangle: 4 curve(s), signed area +0.640000, reversed -0.640000
circle   : 1 curve(s), signed area +0.196032, reversed -0.196032
triangle : 3 curve(s), signed area +0.150000, reversed -0.150000
circle loop: exact rational curve of degree 2, 64 polyline points (16 per span)
untrimmed: is_untrimmed=True, area 4.0000000000
with hole: area 3.2146018366 (4 - pi/4 = 3.2146018366)
contains centre False, contains (0.1, 0.1) True
outer signed area +0.6400, hole -0.1960, area 1.7746018366 (4 * 0.64 - pi/4 = 1.7746018366)
3D loop 0 (outer): 4 curve(s), length 8.00000000
3D loop 1 (hole): 1 curve(s), length 3.14159265
hole circumference 2 pi 0.5 = 3.14159265
normal +1 -> flipped -1, label 'cover', area kept 3.2146018366
quarter of the cylinder, middle half: area 1.5707963268 (2 pi / 4 * 1 = 1.5707963268)
\end{lstlisting}

Points to note:

\begin{itemize}
\item The frame face is created with a \emph{clockwise} outer loop
(\code{rect.reversed()}); the constructor reverses it, and \code{frame.outer.signed_area()}
is positive again.
\item Areas are exact to the integration tolerance: $4-\pi/4$ for the plate with a hole of
radius $0.5$ agrees in all printed digits.
\item On the cylinder (\code{extrude} of a circle: $u$ runs around, $v$ along the axis), the
rectangle $[0, 0.25]\times[0.25, 0.75]$ is a quarter turn of the middle half: area $\pi/2$.
Because the parametrisation of the rational circle is not proportional to the angle, $u=0.25$
is exactly a quarter turn only by the symmetry of the circle's four spans; in general, use
intersections or closest points to find the parameters of a 3D position.
\end{itemize}

\subsection{Area and integrals}\label{sec:trim-area}

\code{area()} uses Green's theorem in the parameter plane. With
$F(u,v) = \int |S_u\times S_v|\,du$, the area is the sum of the loop integrals
$\oint F\,dv$ over all loops; holes, being clockwise, subtract automatically. Both integrals
use a globally adaptive Gauss--Kronrod integrator with a bounded interval budget. The
parameter \code{rel_tol} is the relative accuracy target (default $10^{-9}$). Because the
trim circle is exact, areas of circular holes and discs are exact to that tolerance -- there is
no polygonal approximation anywhere.

\subsection{3D boundaries}\label{sec:trim-boundary}

\code{boundary_curves(tolerance=1e-7)} returns one list of 3D curves per loop (outer loop
first, then the holes). Each trim curve is composed with the surface and fitted by cubic
interpolation with adaptive refinement until it lies within \code{tolerance} of the exact
image $S(c(t))$. In the example, the 3D image of the hole circle has the length
$3.14159265 = 2\pi\cdot0.5$.

\begin{tip}
Where the image of a trim curve is itself exactly representable -- straight segments on planar
faces, iso-parameter lines such as the end circles of a cylinder -- the shell construction and
the STEP writer use the exact curves. \code{boundary_curves} always returns the fitted
approximation; pass a smaller \code{tolerance} when you need more accuracy.
\end{tip}

\subsection{Orientation of a face}\label{sec:trim-flip}

The outward direction of a face is the surface normal $S_u\times S_v$. \code{flipped()}
returns the same face with the opposite normal: the surface is reparametrised with $u$
reversed, $S'(u,v) = S(a+b-u, v)$, and the loops are mirrored so that they still describe
the same region. Area, labels and the 3D boundary are unchanged. You rarely need to flip
faces by hand: \code{Shell.oriented()} does it for all faces of a shell
(\cref{sec:shell-orientation}).

\subsection{Labels}\label{sec:trim-labels}

Every face carries a free-form \code{label} (empty by default), e.g. the name of a boundary
condition or a physical group. \code{with_label(text)} returns a labelled copy. Labels are
kept by \code{flipped}, \code{with_hole}, \code{split}, Boolean operations (every piece of a
split face keeps the label of its face), \code{Shell.oriented} and \code{transformed}, and
they are written to and read from STEP as the face name. The solid primitives label their
faces automatically; see \cref{sec:face-labels} for working with labels on shells.

\section{Splitting surfaces}\label{sec:split}

\apiref{split} cuts a surface into trimmed regions. It has two forms:

\begin{lstlisting}
nf.split(surface, cutter, tolerance=nf.Tolerance())   # along the intersection with a surface
nf.split(surface, pcurves, tolerance=nf.Tolerance())  # along parameter-space curves
\end{lstlisting}

Both return a list of \code{TrimmedSurface} objects that together cover the original surface
without overlap.

\subsection{Splitting by a cutting surface}

\code{split(surface, cutter)} computes the surface/surface intersection
(\cref{sec:isect-ss}) and uses the p-curves on \code{surface} as cuts. Two kinds of cuts
occur:

\begin{itemize}
\item \textbf{Open curves} run from boundary to boundary of the parameter domain (seams count
as boundary). They act as chords of the parameter rectangle: $n$ non-crossing chords give
$n+1$ regions.
\item \textbf{Closed curves} become holes of the region around them and outer loops of new
regions inside -- also nested (a ring between two concentric cuts).
\end{itemize}

If the surfaces do not intersect, the result is a single, untrimmed face.

\subsection{Splitting by parameter-space curves}

\code{split(surface, pcurves)} applies the same rules to curves you provide in the
$(u,v,0)$ convention. Open curves must start and end exactly on the domain boundary
(a curve that overshoots or stops short raises
\code{GeometryError: split: curve ends inside the parameter domain (dangling curve)}) and
should not cross each other.

\begin{warning}
For \code{split(surface, pcurves)} use non-crossing chords. In our tests two \emph{crossing}
chords (a ``+'' through the square) did not produce four quarters but the unsplit face. The
Boolean operations internally use a general planar arrangement that handles crossing cuts
(\cref{sec:boolean-algorithm}), but that variant is not exposed in Python.
\end{warning}

\example{10_split.py}

\begin{lstlisting}[style=text]
plane piece: 1 loop(s), area 10.55575169
plane piece: 2 loop(s), area 25.44424831
disk pi (4 - 0.64) = 10.55575132, frame 36 - disk = 25.44424868
sphere pieces: 2, areas 15.07964474 + 35.18583774; cap 2 pi R h = 15.07964474
no intersection: 1 piece, untrimmed=True
two chords: 3 regions, areas 0.2500, 0.2500, 0.5000
closed curve: 2 regions, areas 0.282743, 0.717257 (pi 0.09 = 0.282743)
\end{lstlisting}

The plane split by the sphere gives the disc inside (one loop) and the frame outside (outer
loop plus a hole). Their areas deviate from the exact values by about $3.7\times10^{-7}$
because the cut is the \emph{fitted} intersection curve (fit tolerance $10^{-7}$ of the model
size), not an exact circle. The sphere split by the plane gives the spherical cap with the
exact area $2\pi R h$ to all printed digits -- there the cut is an iso-parameter circle of the
revolved surface.

\section{Practical advice}\label{sec:trim-advice}

\begin{itemize}
\item \textbf{Map 3D shapes to parameters deliberately.} For a rectangle made with
\code{make_rectangle(origin, edge_u, edge_v)}, $u$ and $v$ are the fractions along the two
edges. A circle of 3D radius $r$ centred at the middle of a square of side $a$ is
\code{circle_loop(0.5, 0.5, r / a)}. For curved surfaces compute loop points with
\code{closest_point} or from intersections.
\item \textbf{Keep loops inside the domain.} Points even slightly outside $[u_0,u_1]\times
[v_0,v_1]$ are rejected. Snap values like $1+10^{-15}$ to the boundary.
\item \textbf{Check results with \code{area()} and \code{contains()}.} They are cheap and
expose wrong orientation or wrong parameter scaling immediately.
\item \textbf{Hand trimmed faces to shells.} A single trimmed face is a surface patch; to get
solids, sew faces with \code{Shell.sew} (\cref{ch:solids}).
\end{itemize}
